\chapter{Anomalous Dimensions of a General Effective Gauge Theory}

%%%%%%%%%%%%%%%%%%%%%%%%%%%
\section{General Effective Gauge Theory}
\label{sec:EFT}
%%%%%%%%%%%%%%%%%%%%%%%%%%%%

In this section, we set the stage for writing down the general EFT. 
We introduce the renormalizable
interactions and the higher-dimensional operators, together with the corresponding notation that is used in subsequent sections.
The complete Lagrangian takes the form
\begin{equation}\label{eq:fullEFT}
    \mathcal{L}_{\text{EFT}} = \mathcal{L}^{(4)}+\mathcal{L}^{(5)}+\mathcal{L}^{(6)}\,, 
\end{equation}
where the individual parts are defined in the following subsections in Eqs.~\eqref{eq:L4}, \eqref{eq:L5}, and \eqref{eq:L6}.


%%%%%%%%%%%%%%%%
\subsection{Renormalizable Interactions}
%%%%%%%%%%%%%%%%

We start by introducing the most general Lagrangian up to mass dimension four that contains an arbitrary number of real scalar fields $\phi_a$, gauge fields with field strengths $F_{\mu\nu}^{A_\alpha}$, and left-handed Weyl spinors $\psi_i$:\footnote{For the two-component notation we follow the conventions in \cite{Dreiner:2008tw}.}
\begin{align}\label{eq:L4}
\mathcal{L}^{(4)} &=  \mathcal{L}^{(4)}_{\text{bos}} + \mathcal{L}^{(4)}_{\text{fer}}\,,\\
    \mathcal{L}^{(4)}_{\text{bos}} &= -\frac{1}{4}\sum_{\alpha=1}^{N_G}\sum_{\beta=1}^{N_G}\xi^{A_\alpha B_\beta} F^{A_\alpha}_{\mu\nu}F^{{B_\beta}\mu\nu}+ \sum_{\alpha=1}^{N_G}\sum_{\beta=1}^{N_G} \tildee \xi^{A_\alpha B_\beta} F^{A_\alpha}_{\mu\nu}\tildee F^{{B_\beta}\,\mu\nu}  + \frac{1}{2} \eta_{ab} (D_\mu\phi)_{a}(D^\mu\phi)_b\nonumber\\
    &\quad -\Lambda-t_a \phi_a-\frac{m^2_{ab}}{2!}\phi_{a}\phi_{b}-\frac{h_{abc}}{3!}\phi_a \phi_b\phi_c
    -\frac{\lambda_{abcd}}{4!}\phi_{a}\phi_{b}\phi_{c}\phi_{d}\,,\\
    \mathcal{L}^{(4)}_{\text{fer}} &= i\, \kappa_{ij}\,  \psi^\dagger_i \overline{\sigma}^\mu (D_\mu \psi)_j - \frac{1}{2!}\left (m_{ij} \psi_i \psi_j + m_{ij}^* \psi_i^\dagger \psi_j^\dagger \right) - \frac{1}{2!}\left (y_{ija} \psi_i \psi_j + y_{ija}^* \psi_i^\dagger \psi_j^\dagger \right)\phi_a  \,.
\end{align}
Here, $\overline \sigma^\mu = (\mathbb{1}_{2\times 2}, -\vec \sigma)$, where $\vec \sigma$ are the Pauli matrices.
%
In addition to all kinetic and mass terms for scalars and fermions, the Lagrangian includes topological terms, the vacuum energy, $\Lambda$, the tadpole interactions, $\{t_a\}$, triple and quartic scalar interactions, $\{h_{abc}\}$ and $\{\lambda_{abcd}\}$, and Yukawa terms, $\{y_{ija}\}$. In particular, the real parameters $m^2_{ab}$, $h_{abc}$, and $\lambda_{abcd}$ are totally symmetric, and the same is true for the complex fermion mass matrix $m_{ij}$ and the Yukawa matrices $y_{ija}$.

The gauge group $G$ is assumed to be a direct product of an arbitrary number of simple and compact gauge groups $G_\alpha$
\begin{equation}
    G = \prod_{\alpha=1}^{N_G} G_\alpha\,.
    \label{eq:full_gauge_group}
\end{equation}
In the following, the adjoint indices of the gauge groups $G_\alpha$ are denoted by $A_\alpha,B_\alpha,\dotsc \in \{1,\dotsc,\dime G_\alpha\}$,
while scalars and fermions are taken to transform in (possibly reducible) representations $S_\alpha$ and $F_\alpha$, respectively. 
The corresponding gauge indices are
$a_{\alpha},b_{\alpha},\dotsc \in \{1,\dotsc,\dime S_{\alpha}\}$ and $i_\alpha,j_\alpha,\dotsc \in \{1,\dotsc,\dime F_{\alpha}\}$, and for notational simplicity their collections over the different gauge-group factors are denoted by $a,b,\dotsc$ and $i,j,\dotsc$, respectively.
Moreover, the covariant derivatives of the scalar and fermion fields are defined by
\begin{align}
    (D_\mu\phi)_{a} &= \partial_\mu\phi_{a} -i\sum_{\alpha=1}^{N_G} g_\alpha A^{A_\alpha}_\mu \theta^{A_\alpha}_{ab}\phi_{b}\,,\label{eq:cov_der_phi} \\
    (D_\mu\psi)_{i} &= \partial_\mu\psi_{i} -i\sum_{\alpha=1}^{N_G} g_\alpha A^{A_\alpha}_\mu t^{A_\alpha}_{ij}\psi_{j}\label{eq:cov_der_psi}\,,
\end{align}
where $g_\alpha$ denotes the gauge couplings and $\theta^{A_\alpha}$ and $t^{A_\alpha}$ are the Hermitian generators acting on the scalar and fermion fields, respectively. They satisfy: 
\begin{gather}
    [\theta^{A_\alpha},\theta^{B_\alpha}] = i\, f^{A_\alpha B_\alpha C_\alpha}\,\theta^{C_\alpha} \,, \quad
    [t^{A_\alpha},t^{B_\alpha}] = i\, f^{A_\alpha B_\alpha C_\alpha}\,t^{C_\alpha} \,,
\end{gather}
%
with the totally anti-symmetric structure constants $f^{A_\alpha B_\alpha C_\alpha}$. The field-strength tensors read\footnote{We use the convention $\epsilon^{0123} =-\epsilon_{0123}= 1$.}
\begin{gather}
    F^{A_\alpha}_{\mu\nu}=\partial_\mu A^{A_\alpha}_\nu - \partial_\nu A^{A_\alpha}_\mu + g_\alpha f^{A_\alpha B_\alpha C_\alpha}A^{B_\alpha}_\mu A^{C_\alpha}_\nu\,,\label{eq:field_strength}\\
    \tildee F^{A_\alpha}_{\mu\nu}= \frac{1}{2}\epsilon_{\mu\nu\rho\sigma}  F^{A_\alpha\rho\sigma}\,.
\end{gather}
%
Finally, due to gauge invariance, the gauge kinetic and topological terms take the following form
\begin{equation}
\xi^{A_\alpha B_\beta} = \delta_{\alpha\beta}\delta^{A_\alpha B_\beta}\,,\quad \tildee \xi^{A_\alpha B_\beta} = \frac{\vartheta_\alpha g_\alpha^2}{32\pi^2} \delta_{\alpha\beta}\delta^{A_\alpha B_\beta} \,,
\end{equation}
with the adjoint indices $A_\alpha$ and $B_\beta$ belonging to non-Abelian gauge factors. 
We normalized the topological angles $\vartheta_\alpha$ in a canonical way using $g^2/(32\pi^2) \int \dd^4x\, F \tildee F \in \mathbb{Z}$.

In the case of Abelian $U(1)$ factors that lead to off-diagonal kinetic and topological terms~\cite{Holdom:1985ag,delAguila:1988jz}, the gauge couplings $g_{\alpha}$ need to be generalized to a matrix $g^{A_\alpha B_\beta}$. In order to allow for this possibility, the simple substitution
\begin{equation}
    g_\alpha \theta^{A_\alpha} \to \sum_{\beta=1}^{N_G} g^{A_\alpha B_\beta} \theta^{B_\beta}
\end{equation} 
can be made in the RGEs below,
and similarly with the fermionic generators.

For convenience, the following definitions are introduced
\begin{align}
    T^\alpha_{abcd} &= \theta^{A_\alpha}_{ab}\theta^{A_\alpha}_{cd}\,,\\
    Y^{A_\alpha B_\beta}_{ab} &= \theta^{A_\alpha}_{ac}\theta^{B_\beta}_{cb}\,,\\
    F^{A_\alpha B_\alpha C_\alpha D_\alpha} &= f^{A_\alpha B_\alpha E_\alpha} f^{C_\alpha D_\alpha E_\alpha}\,,
\end{align}
which occur in several of the RGEs reported in Section~\ref{sec:gEFTresults}.

\subsubsection{Collinear Anomalous Dimensions}

The infrared collinear anomalous dimensions, relevant for the multiplicative renormalization factors, are computed using the on-shell method reviewed in Chapter~\ref{ch:RGEmethod} and exploiting the fact that the stress-energy tensor does not renormalize multiplicatively, as done in \cite{Caron-Huot:2016cwu,EliasMiro:2020tdv,AccettulliHuber:2021uoa,Bresciani:2024shu,Baratella:2022nog,Aebischer:2025zxg}.
They are given by
\begin{align}
    \gamma_{c,v}^{A_\alpha B_\alpha} &= -g_\alpha^2 \left[b_{0,\alpha}\right]^{A_\alpha B_\alpha} \,, \label{eq:coll_vector}  \\
    \gamma_{c,s}^{ab} &= -4\sum_{\alpha=1}^{N_G} g_\alpha^2 \left[C_2(S_\alpha)\right]_{a_{\alpha}b_{\alpha}} + \frac{1}{2}\left( y_{ija}y_{jib}^*+y_{ijb}y_{jia}^* \right) \,, \label{eq:coll_scalar} \\
    \gamma_{c,f}^{ij} &= -3\sum_{\alpha=1}^{N_G} [C_2(F_\alpha)]_{i_\alpha j_\alpha}+\frac{1}{2}y_{ika}y^*_{kja} \,, \label{eq:coll_fermion}
\end{align}
for gauge boson, scalar, and fermion fields, respectively,
with the one-loop $\beta$-function coefficient
\begin{equation}
    \left[b_{0,\alpha}\right]^{A_\alpha B_\alpha} = \frac{11}{3}[C_2(G_\alpha)]^{A_\alpha B_\alpha} - \frac{1}{6} [S_2(S_\alpha)]^{A_\alpha B_\alpha} -\frac{2}{3}[S_2(F_\alpha)]^{A_\alpha B_\alpha}\,.
\end{equation}
The quadratic Casimir $C_2$ and Dynkin indices $S_2$ for the different scalar and fermion representations, $S_\alpha$ and $F_\alpha$, are given by
\begin{align}
    \left[C_2(G_\alpha)\right]^{A_\alpha B_\alpha} &= f^{A_\alpha C_\alpha D_\alpha}f^{B_\alpha C_\alpha D_\alpha}\,,\\
    \left[S_2(S_\alpha)\right]^{A_\alpha B_\alpha} &= \Tr(\theta^{A_\alpha}\theta^{B_\alpha})\,, \\ 
    \left[C_2(S_\alpha)\right]_{a_{\alpha}b_{\alpha}} &= \theta^{A_\alpha}_{a_{\alpha}c_{\alpha}} \theta^{A_\alpha}_{c_{\alpha}b_{\alpha}}\,,\\
    \left[S_2(F_\alpha)\right]^{A_\alpha B_\alpha} &= \Tr(t^{A_\alpha}t^{B_\alpha})\,, \\
    \left[C_2(F_\alpha)\right]_{i_{\alpha}j_{\alpha}} &= t^{A_\alpha}_{j_{\alpha}k_{\alpha}} t^{A_\alpha}_{k_{\alpha}i_{\alpha}}\,.
\end{align}

\subsection{Higher-Dimensional Operator Classification}

A basis of physical operators at a given mass dimension can be conveniently constructed by analyzing the independent kinematic structures associated with the contact amplitudes of the EFT~\cite{Shadmi:2018xan,Cheung:2016drk,Falkowski:2019zdo,Durieux:2019siw,Li:2022tec,DeAngelis:2022qco,Durieux:2019eor}.
In the following, we perform such a construction for operators up to mass dimension six.

We work with the massless spinor-helicity formalism, where a light-like momentum $p$ is decomposed in terms of the helicity spinors $\lambda$ and $\tildee \lambda$ as
\begin{equation}
    p_{\alpha\dot\alpha} = p_\mu \sigma^\mu_{\alpha\dot\alpha} = \lambda_\alpha \tildee \lambda_{\dot\alpha} \,.
\end{equation}
Here, $\sigma^\mu = (\mathbb 1_{2\times 2},\vec \sigma)$ and the spinors $\lambda_\alpha$ and $\tildee \lambda^{\dot\alpha}$ transform in the $(1/2,0)$ and $(0,1/2)$ representations of $SL(2,\mathbb C)$, respectively.
In this framework, the Lorentz-invariant inner products are given by\footnote{We use the convention $\epsilon^{12}=\epsilon^{\dot 1 \dot 2} = -\epsilon_{12}=-\epsilon_{\dot 1 \dot 2} = 1$.}
\begin{gather}
    \agl{i}{j} = \lambda_i{}^\alpha \lambda_j{}_\alpha = \epsilon_{\alpha\beta} \lambda_i{}^\alpha \lambda_j{}^\beta = -\agl{j}{i}\,,\\
    \sqr{i}{j} = \tildee \lambda_i{}_{\dot\alpha} \tildee \lambda_j{}^{\dot \alpha} = -\epsilon_{\dot\alpha\dot\beta} \tildee \lambda_i{}^{\dot\alpha} \tildee \lambda_j{}^{\dot \beta} = -\sqr{j}{i}\,.
\end{gather}
The Mandelstam invariants can then be expressed as 
\begin{equation}
    s_{ij}=2p_i\cdot p_j = \agl{i}{j}\sqr{j}{i}\,.
\end{equation}
Additionally, for real momenta, square and angle products are related via 
\begin{equation}
    \agl{i}{j} = \sqr{j}{i}^*\,.
\end{equation}
Such angle and square brackets are very useful when decomposing scattering amplitudes. In four spacetime dimensions, $n$-point amplitudes have mass dimension $[\mathcal M_n]=4-n$. Furthermore, contact amplitudes $\mathcal{M}$ are polynomials in the inner products and can be schematically written as
\begin{equation}
    \mathcal M = C_i \times \mathcal K_i\,,
\end{equation}
where $C_i$ denotes the Wilson coefficient associated with the operator $\mathcal O_i$, and $\mathcal K_i$ is a kinematic polynomial, $\mathcal K_i = \mathcal K_i(\{\agl{i}{j},\sqr{i}{j}\})$. The mass dimension of $\mathcal K_i$ is fixed by the properties of the operator $\mathcal O_i$:
\begin{equation}
    [\mathcal K_i] = [\mathcal O_i] - \ell(\mathcal O_i) \ge 0\,.
\end{equation}
Here, $\ell(\mathcal O_i)$ denotes the length of the operator, namely the number of fundamental fields it is composed of.
Different kinematic structures can be related via momentum conservation, $\sum_{i=1}^n p_i = 0$, or Schouten identities,
\begin{align}
    \agl{i}{j}\agl{k}{\ell} + \agl{k}{i}\agl{j}{\ell} + \agl{j}{k}\agl{i}{\ell} &= 0\,,\label{eq:schhoutenid1}\\
    \sqr{i}{j}\sqr{k}{\ell} + \sqr{k}{i}\sqr{j}{\ell} + \sqr{j}{k}\sqr{i}{\ell} &= 0\label{eq:schhoutenid2}\,.
\end{align}
Moreover, on-shell three-particle amplitudes vanish due to momentum conservation.
Nevertheless, they can be constructed by assuming complex momenta~\cite{Witten:2003nn}, and, once the particle helicities $\{h_1,h_2,h_3\}$ are fixed, they are uniquely determined by Poincaré invariance, locality, and dimensional analysis~\cite{Benincasa:2007xk}:
\begin{equation}
    \mathcal M(1^{h_1},2^{h_2},3^{h_3}) = 
    \begin{cases}
        C_{\text{H}}\, \agl{1}{2}^{a_3} \agl{2}{3}^{a_1} \agl{3}{1}^{a_2} & \text{if }h_1+h_2+h_3<0\,,\\  C_{\text{A}}\, \sqr{1}{2}^{-a_3} \sqr{2}{3}^{-a_1} \sqr{3}{1}^{-a_2} & \text{if }h_1+h_2+h_3>0\,,
    \end{cases}
\end{equation}
with
\begin{equation}
    a_1 = h_1-h_2-h_3\,,\qquad
    a_2 = h_2-h_3-h_1\,,\qquad
    a_3 = h_3-h_1-h_2\,,
\end{equation}
and $C_{\text{H}}$ and $C_{\text{A}}$ are arbitrary coefficients with mass dimension given by
\begin{equation}
    [C_{\text{H}}] = 1+h_1+h_2+h_3\,,\qquad 
    [C_{\text{A}}] = 1-h_1-h_2-h_3\,.
\end{equation}

Therefore, we can classify all \emph{on-shell} independent kinematic polynomials and the corresponding \emph{physical} operators as follows:
\begin{itemize}
    \item For dimension-five operators, where we have $[\mathcal K_i] = 5 - \ell(\mathcal O_i)$:
    \begin{align}
        \ell_i &= 5
        &&
        \implies 
        &
        \mathcal K_i &=
        \begin{cases}
            1
        \end{cases} 
        &&
        \implies 
        &
        \mathcal O_i &=
        \begin{cases}
            \phi_1\, \phi_2\, \phi_3\, \phi_4\, \phi_5\,,
        \end{cases}
        \\
        \ell_i &= 4 
        &&
        \implies
        &
        \mathcal K_i &= 
        \begin{cases}
             \agl{1}{2}
        \end{cases}
        &&
        \implies 
        &
        \mathcal O_i &=
        \begin{cases}
            \psi_1{}_L\,  \psi_2{}_L\,  \phi_3\,  \phi_4  \,,
        \end{cases}
        \\
        \ell_i &= 3 
        &&
        \implies
        &
        \mathcal K_i &=
        \begin{cases}
            \agl{1}{2}^2 \\
            \agl{1}{2}\agl{1}{3}
        \end{cases}
        &&
        \implies 
        &
        \mathcal O_i &= 
        \begin{cases}
            F_1{}_L\,  F_2{}_L\,  \phi_3\,, \\
            F_1{}_L\,  \psi_2{}_L\,  \psi_3{}_L \,.
        \end{cases}
    \end{align}
    
    \item For dimension-six operators $[\mathcal K_i] = 6 -\ell(\mathcal O_i)$:
    \begin{align}
        \ell_i &= 6
        &&
        \implies 
        &
        \mathcal K_i &=
        \begin{cases}
            1
        \end{cases} 
        &&
        \implies 
        &
        \mathcal O_i &=
        \begin{cases}
            \phi_1\, \phi_2\, \phi_3\, \phi_4\, \phi_5\, \phi_6\,,
        \end{cases}
         \\
        \ell_i &= 5
        &&
        \implies 
        &
        \mathcal K_i &=
        \begin{cases}
            \agl{1}{2}
        \end{cases}
        &&
        \implies 
        &
        \mathcal O_i &=
        \begin{cases}
            \psi_1{}_L\,  \psi_2{}_L\,  \phi_3\,  \phi_4\,  \phi_5 \,,
        \end{cases}
        \\
        \ell_i &= 4
        &&
        \implies 
        &
        \mathcal K_i &= 
        \begin{cases}
            \agl{1}{2}^2 \\
            \agl{1}{2}\agl{1}{3}\\
            \agl{1}{2}\agl{3}{4}\\
            \agl{1}{2}\sqr{3}{4}\\
            \agl{1}{2}\sqr{2}{3} \\
            \agl{1}{2}\sqr{1}{2}
        \end{cases}
        &&
        \implies 
        &
        \mathcal O_i &= 
        \begin{cases}
            F_1{}_L\,  F_2{}_L\,  \phi_3\,  \phi_4\,, \\
            F_1{}_L\,  \psi_2{}_L\,  \psi_3{}_L\,  \phi_4 \,,\\
            \psi_1{}_L\,  \psi_2{}_L\, \psi_3{}_L\,  \psi_4{}_L\,,\\
            \psi_1{}_L\,  \psi_2{}_L\, \psi_3{}_R\,  \psi_4{}_R\,,\\
            \psi_1{}_L\,  D\phi_2\,  \psi_3{}_R\,  \phi_4\,,\\
            D\phi_1\,  D\phi_2\,  \phi_3\,  \phi_4\,.
        \end{cases}
        \\
        \ell_i &= 3
        &&
        \implies 
        &
        \mathcal K_i &=
        \begin{cases}
            \agl{1}{2}\agl{2}{3}\agl{1}{3}
        \end{cases}
        &&
        \implies 
        &
        \mathcal O_i &=
        \begin{cases}
            F_1{}_L\,  F_2{}_L\,  F_3{}_L\,.
        \end{cases}
    \end{align}
\end{itemize}
The scalar and fermion fields are denoted with $\phi$ and $\psi$, respectively, while we use $F$ for the gauge field strengths. The kinematic polynomials with angle brackets substituted by square ones correspond to operators with right-handed fields. 

Finally, the symmetries of the Wilson coefficients are inherited by those of the kinematic structures.
This implies that, for bosonic operators, if a kinematic polynomial is (anti-)symmetric under the exchange of two spinors, then the Wilson coefficient must also be (anti-)symmetric in the associated pair of gauge indices.
For fermionic operators, instead, Fermi-Dirac statistics requires the amplitude to be antisymmetric under the exchange of two identical fermions, so that the above conclusion is reversed: a kinematic polynomial that is symmetric (antisymmetric) under the exchange of two spinors must be contracted with a Wilson coefficient that is antisymmetric (symmetric) in the associated pair of gauge indices.

In the following, we report the basis of operators adopted in this thesis.
Regarding the bosonic operators, we choose to work with Hermitian operators so that all the components of the Wilson coefficients are real.


\subsubsection{Dimension-Five Operators}

The dimension-five Lagrangian can again be split into a bosonic and a fermionic part. It takes the form
%\begingroup
%\interdisplaylinepenalty=10000
\begin{align}
    \mathcal{L}^{(5)} &= \mathcal{L}^{(5)}_{\text{bos}} + \mathcal{L}^{(5)}_{\text{fer}}\,,\label{eq:L5}\\
    \mathcal{L}^{(5)}_{\text{bos}}&=\sum_{\alpha=1}^{N_G}\sum_{\beta=1}^\alpha \left(\left[C_{\phi F^2}\right]^{A_\alpha B_\beta}_{a}\phi_{a}F^{A_\alpha}_{\mu\nu}F^{B_\beta  \mu\nu}+\left[C_{\phi \tildee F^2}\right]^{A_\alpha B_\beta}_{a}\phi_{a}F^{A_\alpha}_{\mu\nu}\tildee F^{B_\beta  \mu\nu}\right)
    \nonumber\\&\quad 
    + \left[C_{\phi^5}\right]_{abcde}\phi_{a}\phi_{b}\phi_{c}\phi_{d}\phi_{e}\,,\\
    \mathcal{L}^{(5)}_{\text{fer}} &= \left[C_{\psi^2\phi^2}\right]_{ijab} \psi_i \psi_j \phi_a \phi_b + \sum_{\alpha=1}^{N_G} \left[C_{\psi^2F}\right]_{ij}^{A_\alpha} \psi_i \sigma^{\mu\nu} \psi_j F_{\mu\nu}^{A_\alpha}  + \text{H.c.}\,,
\end{align}
%\endgroup\ignorespacesafterend
where $\sigma^{\mu\nu} = \frac{i}{4}(\sigma^\mu \overline \sigma^\nu - \sigma^\nu \overline \sigma^\mu)$. In particular, the fermionic operators exhibit non-trivial symmetry properties due to various Dirac relations, which are equivalently reflected in the corresponding kinematic structures. For instance, for anti-commuting spinor fields, one finds:
%
\begin{gather}
    \psi_i \psi_j = \psi_j\psi_i\,, \quad
    \psi_i \sigma^{\mu\nu} \psi_j = -\psi_j \sigma^{\mu\nu}\psi_i\,,
\end{gather}
and analogously for the right-handed spinors. The symmetries of the dimension-five operators together with their minimal form factors are given in Table~\ref{tab:dim5}.

%%%
% --- START TABLE
%%%
\begin{table}[tbp]
\renewcommand{\arraystretch}{2.3}
%\small
\begin{align*}
\resizebox{\textwidth}{!}{
\begin{array}[t]{cccc}
\toprule
\rowcolor{hdrcolor}
\text{\textbf{Name}} & \text{\textbf{Operator}} &  \text{\textbf{Symmetry}} & \text{\textbf{Minimal form factor}} \\
\midrule
\mathcal{O}_{\phi^5} & \phi_a\phi_b\phi_c\phi_d\phi_e & [C_{\phi^5}]_{abcde}=[C_{\phi^5}]_{(abcde)} & F_{\phi^5}(1_{a},2_{b},3_{c},4_{d},5_{e}) = 5! \left[C_{\phi^5}\right]_{abcde}\\
\mathcal{O}_{\phi F^2} & \phi_{a}F^{A_\alpha}_{\mu\nu}F^{B_\beta  \mu\nu} & \left[C_{\phi F^2}\right]^{A_\alpha B_\beta}_{a}=\left[C_{\phi F^2}\right]^{(A_\alpha B_\beta)}_{a} &
F_{\phi F^2}(1_{a},2^-_{A_\alpha},3^-_{B_\beta}) = - \mathcal S_{\alpha\beta} \left[C_{\phi F^2}\right]^{A_\alpha B_\beta}_{a}\agl{2}{3}^2\\
%
\mathcal{O}_{\phi \tildee F^2} & \phi_{a}F^{A_\alpha}_{\mu\nu}\tildee F^{B_\beta  \mu\nu} & \left[C_{\phi \tildee F^2}\right]^{A_\alpha B_\beta}_{a}=\left[C_{\phi \tildee F^2}\right]^{(A_\alpha B_\beta)}_{a} 
 & F_{\phi \tildee F^2}(1_{a},2^-_{A_\alpha},3^-_{B_\beta})  = - i\mathcal S_{\alpha\beta} \left[C_{\phi \tildee F^2}\right]^{A_\alpha B_\beta}_{a}\agl{2}{3}^2\\ 
 %
 \mathcal{O}_{\psi^2\phi^2} & \psi_i \psi_j \phi_a \phi_b  & \left[C_{\psi^2\phi^2}\right]_{ijab} = \left[C_{\psi^2\phi^2}\right]_{(ij)ab}  = \left[C_{\psi^2\phi^2}\right]_{ij(ab)} 
 & F_{\psi^2\phi^2}(1^-_{i},2^-_{j},3_a,4_b)  = 4 \left[C_{\psi^2\phi^2}\right]_{ijab}\agl{1}{2}\\
  %
 \mathcal{O}_{\psi^2 F} &  \psi_i \sigma^{\mu\nu} \psi_j F_{\mu\nu}^{A_\alpha} & \left[C_{\psi^2 F}\right]_{ij}^{A_\alpha} = \left[C_{\psi^2 F}\right]_{[ij]}^{A_\alpha} 
 & F_{\psi^2 F}(1^-_{i},2^-_{j},3^-_{A_\alpha})  = 2\sqrt{2}\left[C_{\psi^2F}\right]_{ij}^{A_\alpha}\agl{1}{3}\agl{2}{3}\\
\bottomrule
%
\end{array}
}
\end{align*}
\caption{Dimension-five bosonic and fermionic operators of the general EFT, together with their minimal form factors and the symmetry properties of the associated Wilson coefficients.}
\label{tab:dim5gen}
\end{table}
%%%
% --- END TABLE
%%%

\subsubsection{Dimension-Six Operators}

The dimension-six interactions are given by the following Lagrangian
%\begingroup
%\interdisplaylinepenalty=10000
\begin{align}\label{eq:L6}
    \mathcal{L}^{(6)} &= \mathcal{L}^{(6)}_{\text{bos}}+\mathcal{L}^{(6)}_{\text{fer}}\,,\\
    \mathcal{L}^{(6)}_{\text{bos}}&=\sum_{\alpha=1}^{N_G}\sum_{\beta=1}^\alpha \left(\left[C_{\phi^2 F^2}\right]^{A_\alpha B_\beta}_{ab}\phi_{a}\phi_{b}F^{A_\alpha}_{\mu\nu}F^{B_\beta  \mu\nu}+\left[C_{\phi^2 \tildee F^2}\right]^{A_\alpha B_\beta}_{ab}\phi_{a}\phi_{b}F^{A_\alpha}_{\mu\nu}\tildee F^{B_\beta  \mu\nu}\right)
    \nonumber\\&\quad  \nonumber
    + \left[C_{\phi^6}\right]_{abcdef}\phi_{a}\phi_{b}\phi_{c}\phi_{d}\phi_{e}\phi_{f}+\left[C_{D^2\phi^4}\right]_{abcd}(D_\mu \phi)_{a}(D^\mu \phi)_{b}\phi_{c}\phi_{d}\\&\quad
    +\sum_{\alpha=1}^{N_G}\left(\left[C_{F^3}\right]^{A_\alpha B_\alpha C_\alpha}F^{A_\alpha \nu}_{\mu}F^{B_\alpha \rho}_{ \nu} F^{C_\alpha \mu}_{\rho}+\left[C_{\widetilde F^3}\right]^{A_\alpha B_\alpha C_\alpha}F^{A_\alpha \nu}_{\mu}F^{B_\alpha\rho}_{ \nu} \widetilde F^{C_\alpha\mu}_{\rho}\right)\,,\\
    \mathcal{L}^{(6)}_{\text{fer}} &= \bigg(\left[C_{\psi^2\phi^3}\right]_{ijabc} \psi_i \psi_j \phi_a \phi_b \phi_c + \sum_{\alpha=1}^{N_G} \left[C_{\psi^2 \phi F}\right]_{ija}^{A_\alpha} \psi_i \sigma^{\mu\nu} \psi_j \phi_a F_{\mu\nu}^{A_\alpha}  \nonumber \\&\quad + \left[C_{\psi^4}\right]_{ijkl} (\psi_i \psi_j)(\psi_k\psi_l) + \text{H.c.}\bigg) + \left[C_{\overline \psi^2 \psi^2}\right]_{ijkl}(\psi_i^\dagger \overline \sigma^\mu \psi_j)(\psi_k^\dagger \overline \sigma_\mu \psi_l)\nonumber \\&\quad
    + i\left[C_{D\overline \psi \psi \phi^2}\right]_{ijab}(\psi^\dagger_i \overline \sigma^\mu \psi_j)\left[(D_\mu \phi)_a \phi_b - \phi_a (D_\mu \phi)_b\right]\,.
\end{align}
%\endgroup\ignorespacesafterend
The symmetries of the corresponding Wilson coefficients are collected for bosonic and fermionic operators in Table~\ref{tab:dim6}, and the minimal form factors are given in Table~\ref{tab:dim6FFs}. 

%%%
% --- START TABLE
%%%
\begin{table}[tbp]
\renewcommand{\arraystretch}{2.3}
%\small
\begin{align*}
\resizebox{\textwidth}{!}{
\begin{array}[t]{cccc}
\toprule
\rowcolor{hdrcolor}
\text{\textbf{Name}} & \text{\textbf{Operator}} &  \text{\textbf{Symmetry}} \\
\midrule
\mathcal O_{\phi^6} & \phi_a\phi_b\phi_c\phi_d\phi_e\phi_f & [C_{\phi^6}]_{abcdef}=[C_{\phi^6}]_{(abcdef)} \\
%%%%%%%%
\mathcal O_{D^2\phi^4}  & (D_\mu \phi)_{a}(D^\mu \phi)_{b}\phi_{c}\phi_{d} & [C_{D^2\phi^4}]_{abcd} = [C_{D^2\phi^4}]_{(ab)cd} = [C_{D^2\phi^4}]_{ab(cd)}
\\
%%%%%%%%
\mathcal{O}_{\phi^2 F^2} & \phi_{a}\phi_{b}F^{A_\alpha}_{\mu\nu}F^{B_\beta  \mu\nu} & \left[C_{\phi^2F^2}\right]^{A_\alpha B_\beta}_{ab}=\left[C_{\phi^2F^2}\right]^{(A_\alpha B_\beta)}_{ab}=\left[C_{\phi^2F^2}\right]^{A_\alpha B_\beta}_{(ab)} \\
%%%%%%%%
\mathcal{O}_{\phi^2\tildee F^2} & \phi_{a}\phi_{b}F^{A_\alpha}_{\mu\nu}\tildee F^{B_\beta  \mu\nu} & \left[C_{\phi^2 \tildee F^2}\right]^{A_\alpha B_\beta}_{ab}=\left[C_{\phi^2\tildee F^2}\right]^{(A_\alpha B_\beta)}_{ab}=\left[C_{\phi^2\tildee F^2}\right]^{A_\alpha B_\beta}_{(ab)}\\
 %%%%%%%%
 \mathcal{O}_{F^3} & F^{A_\alpha\nu}_{\mu}F^{B_\alpha\rho}_{ \nu}F^{C_\alpha\mu}_{\rho} & \left[C_{F^3}\right]^{A_\alpha B_\alpha C_\alpha}=\left[C_{F^3}\right]^{[A_\alpha B_\alpha C_\alpha]}\\
 %%%%%%%%
  \mathcal{O}_{\tildee F^3} & F^{A_\alpha\nu}_{\mu}F^{B_\alpha\rho}_{ \nu}{\tildee F}^{C_\alpha\mu}_{\rho} & \left[C_{\tildee F^3}\right]^{A_\alpha B_\alpha C_\alpha}=\left[C_{\tildee F^3}\right]^{[A_\alpha B_\alpha C_\alpha]}\\
 %%%%%%%%
%\midrule
\mathcal O_{\psi^2\phi^3} &  \psi_i \psi_j \phi_a \phi_b \phi_c & \left[C_{\psi^2\phi^3}\right]_{ijabc} = \left[C_{\psi^2\phi^3}\right]_{ij(abc)} = \left[C_{\psi^2\phi^3}\right]_{(ij)abc} \\
%%%%%%%%
\mathcal O_{\psi^2 \phi F}  & \psi_i \sigma^{\mu\nu} \psi_j \phi_a F_{\mu\nu}^{A_\alpha} & \left[C_{\psi^2 \phi F}\right]_{ija}^{A_\alpha} = \left[C_{\psi^2 \phi F}\right]_{[ij]a}^{A_\alpha}
\\
%%%%%%%%
\mathcal{O}_{\psi^4} & (\psi_i \psi_j)(\psi_k\psi_l) & \left[C_{\psi^4}\right]_{ijkl} = \left[C_{\psi^4}\right]_{(ij)kl} = \left[C_{\psi^4}\right]_{ij(kl)} = \left[C_{\psi^4}\right]_{klij}  \\
%%%%%%%%
\mathcal{O}_{\overline \psi^2 \psi^2} & (\psi_i^\dagger \overline \sigma^\mu \psi_j)(\psi_k^\dagger \overline \sigma_\mu \psi_l) & \left[C_{\overline \psi^2 \psi^2}\right]_{ijkl} = \left[C_{\overline \psi^2 \psi^2}\right]_{kjil} = \left[C_{\overline \psi^2 \psi^2}\right]_{ilkj} = \left[C_{\overline \psi^2 \psi^2}\right]_{jilk}^*\\
 %%%%%%%%
 \mathcal{O}_{D\overline \psi \psi \phi^2} & i(\psi^\dagger_i \overline \sigma^\mu \psi_j)\left[(D_\mu \phi)_a \phi_b - \phi_a (D_\mu \phi)_b\right] & \left[C_{D\overline \psi \psi \phi^2}\right]_{ijab} = \left[C_{D\overline \psi \psi \phi^2}\right]_{ij[ab]} = \left[C_{D\overline \psi \psi \phi^2}\right]_{jiba}^*\\
 %%%%%%%%
\bottomrule
%
\end{array}
}
\end{align*}
\caption{Dimension-six bosonic and fermionic operators of the general EFT, together with the symmetry properties of the associated Wilson coefficients.}
\label{tab:dim6gen}
\end{table}
%%%
% --- END TABLE
%%%
%%%
% --- START TABLE
%%%
\begin{table}[tbp]
\renewcommand{\arraystretch}{2.3}
%\small
\begin{align*}
%\resizebox{\textwidth}{!}{
\begin{array}[t]{cc}
\toprule
\rowcolor{hdrcolor}
\text{\textbf{Name}}  & \text{\textbf{Minimal form factor}} \\
\midrule
\mathcal O_{\phi^6} & F_{\phi^6}(1_{a},2_{b},3_{c},4_{d},5_{e},6_{f}) = 6! \left[C_{\phi^6}\right]_{abcdef}\\
%%%%%%%%
\mathcal O_{D^2\phi^4}  & F_{D^2\phi^4}(1_a,2_b,3_c,4_d) =-2\left(\left[\widehat C_{D^2\phi^4}\right]_{abcd} \agl{1}{2}\sqr{2}{1}  + \left[\widehat C_{D^2\phi^4}\right]_{acbd} \agl{1}{3}\sqr{3}{1}\right)\\
%%%%%%%%
\mathcal{O}_{\phi^2 F^2} &
F_{\phi^2 F^2}(1_a,2_b,3^-_{A_\alpha},4^-_{B_\beta}) = -2 \mathcal S_{\alpha\beta} \left[C_{\phi^2 F^2}\right]^{A_\alpha B_\beta}_{a b}\agl{3}{4}^2\\
%%%%%%%%
\mathcal{O}_{\phi^2\tildee F^2} & F_{\phi^2 \tildee F^2}(1_a,2_b,3^-_{A_\alpha},4^-_{B_\beta}) = -2i \mathcal S_{\alpha\beta} \left[C_{\phi^2 \tildee F^2}\right]^{A_\alpha B_\beta}_{a b}\agl{3}{4}^2\\
 %%%%%%%%
 \mathcal{O}_{F^3} & F_{F^3}(1^-_{A_\alpha},2^-_{B_\alpha},3^-_{C_\alpha}) =-3i\sqrt{2} \left[C_{F^3}\right]^{A_\alpha B_\alpha C_\alpha} \agl{1}{2}\agl{2}{3}\agl{3}{1}\\
 %%%%%%%%
  \mathcal{O}_{\tildee F^3} & F_{\tildee F^3}(1^-_{A_\alpha},2^-_{B_\alpha},3^-_{C_\alpha}) =3\sqrt{2} \left[C_{\tildee F^3}\right]^{A_\alpha B_\alpha C_\alpha} \agl{1}{2}\agl{2}{3}\agl{3}{1}\\
 %%%%%%%%
% \midrule
\mathcal O_{\psi^2\phi^3} & F_{\psi^2 \phi^3}(1^-_i,2^-_j,3_a,4_b,5_c) = 12 \left[ C_{\psi^2\phi^3}\right]_{ijabc} \agl{1}{2}\\
%%%%%%%%
\mathcal O_{\psi^2 \phi F}   & F_{\psi^2 \phi F}(1^-_i,2^-_j,3_a,4^-_{A_\alpha}) = 2\sqrt{2} \left[C_{\psi^2 \phi F}\right]_{ija}^{A_\alpha}\agl{1}{4}\agl{2}{4}\\
%%%%%%%%
\mathcal{O}_{\psi^4} &
F_{\psi^4}(1^-_i,2^-_j,3^-_k,4^-_l) = 8 \left( \left[\widehat C_{\psi^4}\right]_{ijkl}\agl{1}{2}\agl{3}{4} + \left[\widehat C_{\psi^4}\right]_{ikjl}\agl{1}{3}\agl{4}{2}\right) \\
%%%%%%%%
\mathcal{O}_{\overline \psi^2 \psi^2} & F_{\overline \psi^2 \psi^2}(1^+_i,2^-_j,3^+_k,4^-_l) = 8 \left[C_{\overline \psi^2 \psi^2}\right]_{ijkl}\agl{2}{4}\sqr{3}{1}\\
 %%%%%%%%
\mathcal{O}_{D\overline \psi \psi \phi^2} & F_{D\overline \psi \psi \phi^2}(1^+_i,2^-_j,3_a,4_b) = -4 \left[C_{D\overline \psi \psi \phi^2}\right]_{ijab}\agl{2}{3}\sqr{3}{1}\\
 %%%%%%%%
\bottomrule
%
\end{array}
%}
\end{align*}
\caption{Dimension-six bosonic and fermionic operators and their minimal form factors. We have defined the following combinations of Wilson coefficients: $[\widehat C_{D^2\phi^4}]_{abcd} = [C_{D^2\phi^4}]_{abcd}+[C_{D^2\phi^4}]_{cdab}-[C_{D^2\phi^4}]_{adbc}-[C_{D^2\phi^4}]_{bcad}$ and $[\widehat C_{\psi^4}]_{ijkl} = [C_{\psi^4}]_{ijkl} - [C_{\psi^4}]_{iljk}$.}
\label{tab:dim6FFs}
\end{table}
%%%
% --- END TABLE
%%%

For instance, the four-fermion operator $\mathcal{O}_{\overline \psi^2 \psi^2}$ satisfies the following Fierz relation:
%
\begin{equation}
    (\psi_i^\dagger \overline \sigma_\mu \psi_j)(\psi_k^\dagger \overline \sigma^\mu \psi_l) = 2 (\psi_i^\dagger \psi_k^\dagger)( \psi_l \psi_j) \,,
\end{equation}
which leads to several symmetry relations. Furthermore, using the identity
%
\begin{equation}
    (\psi_i^\dagger \overline \sigma_\mu \psi_j) = - (\psi_j \sigma_\mu \psi_i^\dagger)\,,
\end{equation}
the operator $\mathcal{O}_{\overline \psi^2 \psi^2}$ describes four-fermion vector operators with both equal and mixed chiralities. 
Moreover, as a consequence of the Schouten identities in Eqs.~\eqref{eq:schhoutenid1} and \eqref{eq:schhoutenid2},
the Wilson coefficient associated with the operator $\mathcal O_{\psi^4}$ satisfies
\begin{equation}
    \left[C_{\psi^4}\right]_{ijkl} + \left[C_{\psi^4}\right]_{iklj} + \left[C_{\psi^4}\right]_{iljk} = 0\,,
\end{equation}
in addition to the identities reported in Table~\ref{tab:dim6}.

Let us further note a subtlety that involves the $D^2\phi^4$-class operator $[\mathcal O_{D^2\phi^4}]_{abcd} = (D_\mu \phi)_a (D^\mu \phi)_b \phi_c \phi_d$. 
For this, we define the following combination of Wilson coefficients:
\begin{equation}\label{eq:hatCD2phi4}
\left[\widehat C_{D^2\phi^4}\right]_{abcd}=\left[C_{D^2\phi^4}\right]_{abcd}+\left[C_{D^2\phi^4}\right]_{cdab}-\left[C_{D^2\phi^4}\right]_{adbc}-\left[C_{D^2\phi^4}\right]_{bcad}\,.
\end{equation}
Adopting a geometric approach, this combination can be shown to be proportional to the Riemann tensor of the scalar metric,\footnote{For a scalar metric $g_{ab}(\phi) = \delta_{ab} + 2[C_{D^2\phi^4}]_{abcd}\phi^c \phi^d$, the corresponding Riemann tensor $R_{abcd}$ is proportional to $[\widehat C_{D^2\phi^4}]_{acbd} = [C_{D^2\phi^4}]_{acbd}+[C_{D^2\phi^4}]_{bdac}-[C_{D^2\phi^4}]_{adbc}-[C_{D^2\phi^4}]_{bcad}$.} which implies the following symmetries
\begin{gather}
    \left[\widehat C_{D^2\phi^4}\right]_{abcd} = \left[\widehat C_{D^2\phi^4}\right]_{badc} = -\left[\widehat C_{D^2\phi^4}\right]_{adcb} = -\left[\widehat C_{D^2\phi^4}\right]_{cbad}\,,\\
    \left[\widehat C_{D^2\phi^4}\right]_{abcd} + \left[\widehat C_{D^2\phi^4}\right]_{acdb} + \left[\widehat C_{D^2\phi^4}\right]_{adbc} = 0\,.
\end{gather}
Furthermore, by introducing the redundant operator
\begin{equation}
\left[\mathcal R_{D^2\phi^4}\right]_{abcd} = (\Box \phi)_a \phi_b \phi_c \phi_d
\end{equation}
and by using integration by parts, one finds the relations
\begin{align}
    \left[\mathcal O_{D^2\phi^4}\right]_{abcd} - \left[\mathcal O_{D^2\phi^4}\right]_{cdab} &= -\frac{1}{2}\left(
    \left[\mathcal R_{D^2\phi^4}\right]_{abcd} + \left[\mathcal R_{D^2\phi^4}\right]_{bacd}\right.\nonumber\\
    &\left.\qquad\,\,-\left[\mathcal R_{D^2\phi^4}\right]_{cabd}-\left[\mathcal R_{D^2\phi^4}\right]_{dabc}\,
    \right)
\end{align}
and
\begin{equation}
    \left[\mathcal O_{D^2\phi^4}\right]_{abcd} + \left[\mathcal O_{D^2\phi^4}\right]_{adbc} +
    \left[\mathcal O_{D^2\phi^4}\right]_{acdb}  = -\left[\mathcal R_{D^2\phi^4}\right]_{abcd}\,,\label{eq:D2phi4symmetric}
\end{equation}
implying
\begin{equation}
\left[\mathcal O_{D^2\phi^4}\right]_{(abcd)} = -\frac{1}{3}\left[\mathcal R_{D^2\phi^4}\right]_{(abcd)}\,.
\end{equation}
One can then show, by using the equation of motion of the scalar field,
\begin{equation}
    (\Box \phi)_a = -t_a -m^2_{ab}\,\phi_b - \frac{1}{2!}h_{abc}\,\phi_b\phi_c - \frac{1}{3!}\lambda_{abcd}\,\phi_b\phi_c\phi_d
    -\frac{1}{2!}\left(y_{ija}\psi_i \psi_j + y_{ija}^* \psi_i^\dagger \psi_j^\dagger
    \right)
    \,,
\end{equation}
that the redundant operator $\mathcal R_{D^2\phi^4}$ is a linear combination of $\mathcal O_{\phi^6}$, $\mathcal O_{\phi^5}$, $\mathcal O_{\phi^4}$, $\mathcal O_{\phi^3}$, and $\mathcal O_{\psi^2\phi^3}$.
Consequently, generating $\mathcal R_{D^2\phi^4}$ at the one-loop level influences the RGEs of the coefficients $C_{\phi^6}$, $C_{\phi^5}$, $\lambda$, $h$, and $C_{\psi^2\phi^3}$.
To illustrate this point, we write the Wilson coefficient $[C_{D^2\phi^4}]_{abcd}$ as
\begin{equation}
    \left[C_{D^2\phi^4}\right]_{abcd} = \frac{1}{6}\left(
    2\left[\widehat C_{D^2\phi^4}\right]_{abcd}-\left[\widehat C_{D^2\phi^4}\right]_{acbd}
    \right)+\left[\tildee C_{D^2\phi^4}\right]_{abcd}+\left[\overline C_{D^2\phi^4}\right]_{abcd}\,,
    \label{eq:decomposition}
\end{equation}
where we introduced the combinations
\begin{align}
    \left[\tildee C_{D^2\phi^4}\right]_{abcd} &= \left[ C_{D^2\phi^4}\right]_{(abcd)}\,,\label{eq:Ctilde_def}\\
    \left[\overline C_{D^2\phi^4}\right]_{abcd} &= \frac{1}{2}\left(
    \left[C_{D^2\phi^4}\right]_{abcd} - \left[C_{D^2\phi^4}\right]_{cdab}
    \right)\,.\label{eq:Cbar_def}
\end{align}
In this notation, one finds the following contributions from $\mathcal R_{D^2\phi^4}$
\begin{align}
    \left[\dot C_{\psi^2\phi^3}\right]_{ijabc} &\supset \frac{1}{3!}\frac{1}{2!}\sum_{\sigma(\{a,b,c\})}y_{ijd} \left(
    \frac{1}{3}\left[\dot{\tildee C}_{D^2\phi^4}\right]_{dabc}+\left[\dot{\overline C}_{D^2\phi^4}\right]_{dabc}
    \right)\,,\label{eq:psi2phi3<-Ctildeandbar}\\
    \left[\dot C_{\phi^6}\right]_{abcdef} &\supset \frac{1}{6!}\frac{1}{3!}\sum_{\sigma(\{a,b,c,d,e,f\})}\lambda_{abcg}\left(
    \frac{1}{3}\left[\dot{\tildee C}_{D^2\phi^4}\right]_{gdef}+\left[\dot{\overline C}_{D^2\phi^4}\right]_{gdef}
    \right)\,,\label{eq:phi6<-Ctildeandbar}\\
    \left[\dot C_{\phi^5}\right]_{abcde} &\supset \frac{1}{5!}\frac{1}{2!}\sum_{\sigma(\{a,b,c,d,e\})}h_{abf}\left(
    \frac{1}{3}\left[\dot{\tildee C}_{D^2\phi^4}\right]_{fcde}+\left[\dot{\overline C}_{D^2\phi^4}\right]_{fcde}
    \right)\,,\label{eq:phi5<-Ctildeandbar}\\
    \dot \lambda_{abcd} &\supset - \sum_{\sigma(\{a,b,c,d\})}m^2_{ae}\left(
    \frac{1}{3}\left[\dot{\tildee C}_{D^2\phi^4}\right]_{ebcd}+\left[\dot{\overline C}_{D^2\phi^4}\right]_{ebcd}
    \right)\,,\label{eq:phi4<-Ctildeandbar}\\
    \dot h_{abc} &\supset -  \sum_{\sigma(\{a,b,c\})} t_d \left(
    \frac{1}{3}\left[\dot{\tildee C}_{D^2\phi^4}\right]_{dabc}+\left[\dot{\overline C}_{D^2\phi^4}\right]_{dabc}
    \right)\,,\label{eq:phi3<-Ctildeandbar}
\end{align}
which have to be taken into account and added to the results presented in Section~\ref{sec:gEFTresults}.
With $\sum_{\sigma(I)}$ we denote the sum over all permutations of the set of indices $I$, and $\sum_{\sigma(I\times J)} = \sum_{\sigma(I)}\sum_{\sigma(J)}$.

\section{Results\label{sec:gEFTresults}}